\chapter{The Baxter platform}
\label{ch:platform}

This chapter is background: what the hardware is, what its joints are called, and
which of the legacy ROS~1 concepts survive into this workspace. Everything here
is factual description of the robot, drawn from the Rethink documentation set and
from what the robot was observed to do during the hardware sessions.

\section{The robot}

Baxter is a two-armed research robot on a fixed pedestal. Each arm has seven
revolute joints; the arms are series-elastic, so every joint has a spring between
the motor and the link. That single design choice explains most of what
Chapter~\ref{ch:hwmotion} measures: the arm is compliant by construction, it does
not track a fast setpoint rigidly, and lag is a physical property rather than a
tuning defect.

The workspace also carries a head with a pan joint, a screen, and two
one-degree-of-freedom electric grippers. This workspace commands the arms only.
The head pan and the two gripper finger joints are published as state and are not
commanded; the model provides them so that TF and the robot display are complete.

\section{Joint naming}

Joints are named \texttt{<side>\_<segment>}, where side is \texttt{left} or
\texttt{right}. The seven suffixes, in chain order from the torso outward:

\begin{center}
\resizebox{\linewidth}{!}{% The seven arm joints in chain order, with the URDF limits the shim enforces.
% Values from follow_joint_trajectory_shim.py JOINT_LIMITS, transcribed from
% baxter_description/urdf/baxter_base/baxter_base.urdf.xacro.
\begin{tikzpicture}[
  font=\small,
  j/.style={draw, circle, minimum size=9mm, inner sep=0pt, fill=notebg,
            draw=noterule, font=\footnotesize\ttfamily},
  seg/.style={-{Stealth[length=2mm]}, thick, draw=coderule},
  lim/.style={font=\scriptsize, align=center, text=comment},
  grp/.style={font=\scriptsize\bfseries, text=noterule},
]

\node[j] (s0) at (0,0) {s0};
\node[j] (s1) at (2.1,0) {s1};
\node[j] (e0) at (4.2,0) {e0};
\node[j] (e1) at (6.3,0) {e1};
\node[j] (w0) at (8.4,0) {w0};
\node[j] (w1) at (10.5,0) {w1};
\node[j] (w2) at (12.6,0) {w2};

\foreach \a/\b in {s0/s1, s1/e0, e0/e1, e1/w0, w0/w1, w1/w2}
  \draw[seg] (\a) -- (\b);

\node[left=3mm of s0, align=right, font=\footnotesize] {torso\\mount};
\node[right=3mm of w2, align=left, font=\footnotesize] {gripper};

\node[lim, below=3mm of s0] {$\pm 1.702$\\1.5\,rad/s};
\node[lim, below=3mm of s1] {$-2.147 \ldots 1.047$\\1.5\,rad/s};
\node[lim, below=3mm of e0] {$\pm 3.054$\\1.5\,rad/s};
\node[lim, below=3mm of e1] {$-0.05 \ldots 2.618$\\1.5\,rad/s};
\node[lim, below=3mm of w0] {$\pm 3.059$\\4.0\,rad/s};
\node[lim, below=3mm of w1] {$-1.571 \ldots 2.094$\\4.0\,rad/s};
\node[lim, below=3mm of w2] {$\pm 3.059$\\4.0\,rad/s};

\node[grp] at (1.05,0.95) {shoulder};
\node[grp] at (5.25,0.95) {elbow};
\node[grp] at (10.5,0.95) {wrist};

\draw[decorate, decoration={brace, amplitude=3pt}, draw=noterule]
  (-0.5,0.7) -- (2.6,0.7);
\draw[decorate, decoration={brace, amplitude=3pt}, draw=noterule]
  (3.7,0.7) -- (6.8,0.7);
\draw[decorate, decoration={brace, amplitude=3pt}, draw=noterule]
  (7.9,0.7) -- (13.1,0.7);

\end{tikzpicture}
}
\end{center}

\begin{itemize}
  \item \texttt{s0}, \texttt{s1} --- shoulder rotation and flexion.
  \item \texttt{e0}, \texttt{e1} --- upper-arm roll and elbow flexion.
  \item \texttt{w0}, \texttt{w1}, \texttt{w2} --- wrist roll, flexion, and final roll.
\end{itemize}

Two features of that table matter later. First, \texttt{s1} is the only strongly
asymmetric joint ($-2.147$ to $1.047$\,rad): a move that fits in one direction may
not fit in the other, which is why the example clients pick a reversible target
rather than a fixed offset. Second, the three wrist joints permit
$4.0$\,rad/s against $1.5$\,rad/s for the shoulder and elbow --- so a planner that
distributes motion evenly across joints will drive the wrists fastest. That is
the mechanism behind the MoveIt hardware limitation in
Section~\ref{sec:moveit-hardware}.

\section{Tucked and untucked}

Baxter parks with its arms folded against the torso. In that pose \texttt{s1}
sits at approximately $-2.175$\,rad, which is \emph{outside} the URDF limit table
this workspace enforces ($-2.147$). This is not a transcription error: the parked
pose is a mechanical rest position rather than a commanded one.

The consequence is deliberate and worth stating early. This workspace's action
shim rejects any goal containing an out-of-limit position, so while the robot is
tucked it will refuse to move --- correctly. Untucking is therefore done with the
legacy Rethink tool, not through the ROS~2 path, and Chapter~\ref{ch:bridge}
explains why replacing that tool is the one part of the legacy stack this project
does not simply reimplement.

\section{Enabling}

Baxter's arms are unpowered until the robot is explicitly enabled, and it reports
its state on \texttt{/robot/state} as an \texttt{AssemblyState}: enabled, stopped,
in error, and whether an estop is engaged. Motion requires enabled, not stopped,
no error.

An enable can be refused. If the robot believes its arms are in collision --- which
the tucked pose can trigger --- the enable is dropped almost immediately.
The legacy untuck tool works around this by suppressing collision avoidance on
both limbs at 20\,Hz \emph{while} it republishes the enable, rather than enabling
once and waiting. That detail is not a curiosity: it is the difference between an
enable that holds and one that does not, and it was rediscovered the hard way
during the 2026-07-22 session.

\begin{safetybox}
An enabled Baxter has powered arms. In gravity-compensation mode --- enabled but
uncommanded --- the arms hold position but move freely if pushed, and they will
drift. Never enable the robot without the e-stop in reach and the workspace clear.
\end{safetybox}

\section{The legacy ROS 1 interface}

The robot runs its own ROS~1 Noetic master. Everything this workspace does on
hardware is a client of that master. The topics that matter:

\begin{table}[h]
\centering
\small
\begin{tabularx}{\linewidth}{@{}llX@{}}
\toprule
\textbf{Topic} & \textbf{Type} & \textbf{Role} \\
\midrule
\texttt{/robot/state} & \texttt{AssemblyState} & Enabled, stopped, error, estop. The safety gate reads this. \\
\texttt{/robot/joint\_states} & \texttt{JointState} & Measured positions. Published by more than one node. \\
\texttt{/robot/ref\_joint\_states} & \texttt{JointState} & The robot's own commanded reference, $\approx$140\,Hz. \\
\texttt{/robot/limb/*/joint\_command} & \texttt{JointCommand} & Position commands in. \\
\texttt{robot/set\_super\_enable} & \texttt{Bool} & Enable/disable. \\
\bottomrule
\end{tabularx}
\end{table}

\begin{notebox}[One measurement worth carrying forward]
\texttt{/robot/joint\_states} is published by \emph{two} nodes --- arm joints from
one, gripper joints from another. A subscriber that keeps only the latest message
therefore sees an incomplete picture roughly a third of the time. Any client that
reads joint positions on hardware has to merge across messages by name. This is
handled in the workspace's clients and is called out again in
Section~\ref{sec:jointstate-merge}, because it is the kind of defect that looks
like flaky hardware.
\end{notebox}

Commanded motion is not fire-and-forget. \texttt{JointCommand} has a timeout: if
commands stop arriving for longer than \texttt{joint\_command\_timeout}
(0.2\,s by default), the arm reverts to gravity compensation. Holding a position
means continuing to command it, which is why the shim in
Chapter~\ref{ch:bridge} publishes at 100\,Hz for as long as a goal is active and
keeps commanding the final point after the trajectory ends.
